\documentclass[../../../include/open-logic-section]{subfiles}

\begin{document}

\olfileid{sth}{story}{extensionality}
\olsection{د غړو له مخې برابري}

تر ټولو لومړنۍ خبره دا ده چې سټونه د خپلو !!{element}s له مخې
له يو بل څخه بېلېږي۔ په لا کره توګه:

\begin{axiom}[د غړو له مخې برابري]
که سټونه $A$ او $B$ يو شان !!{element}s ولري، نو $A$ او $B$
هماغه يو سټ دے۔
\[
  \lforall[A][\lforall[B][(\lforall[x][(x \in A \liff x \in B)] \lif
  \eq[A][B])]]
\]
\end{axiom}

دا مو د \olref[sfr][][]{part} په ټوله برخه کښې فرض کړې وه۔
خو بيا ويل يې ګټه لري۔ د غړو له مخې د برابرۍ اصل دا بنسټيز
تصور بيانوي چې يو سټ د خپلو !!{element}s له مخې ټاکل کېږي۔
(نو سټونه له \emph{مفهومونو} سره پرتله کولے شو، چې عين هماغه
شيان ښايي د ډېرو بېلو مفهومونو لاندې راشي۔)

ولې دا اصل ومنو؟ دا خبره د منلو وړ ده چې د غړو له مخې د
برابرۍ هر انکار د داسې هر څه د پرېښودو پرېکړه ده چې ورته
\emph{سټ تيوري} ويل کېدے شي۔ سټ تيوري، نۀ زياته نۀ کمه،
د هغو ټولګو تيوري ده چې د خپلو غړو له مخې ټاکل کېږي۔

خو په \olref[sth][][]{part} کښې اصلي ننګونه د هغو اصلونو
ټاکل دي چې راته ووايي \emph{کوم سټونه وجود لري}۔ او څرګندېږي
چې د دې پوښتنې يوازې هغه ځواب چې په رښتيا «څرګند» برېښي،
په ثبوت سره ناسم دے۔

\end{document}
